\documentclass[book.tex]{subfiles}

\begin{document}

\chapter{The Dirac Equation}

\label{chap:dirac_equation}
\noindent\rule{11cm}{0.4pt} \\
\textbf{Prerequisites:}  \autoref{chap:schrodinger} \\
\textbf{Difficulty Level:} ** \\
\noindent\rule{11cm}{0.4pt}
\vspace{5mm}

The Dirac equation provides a mechanism for modeling relativistic quantum mechanical systems. Wave functions as used in the Dirac equation are bispinors instead of the usual complex numbers seen elsewhere in quantum mechanics.

\section{A Brief Introduction to Spinors}

The use of spinors is needed to model the intrinsic spin of the system. The wave-function on spacetime is consequently of the form

\begin{align}
    \psi: \mathbb{R}^{1, 3} \to \mathbb{C}^4
\end{align}

The space these functions live in is given by 
\begin{align}
    \mathcal{H} &= L^2(\mathbb{R}^{1, 3}, \mathbb{C}^4)
\end{align}

The bispinor is given by vectors in $\mathbb{C}^4$. The Dirac equation itself is written in terms of a bispinor field 
\begin{align}
    \mathbb{R}^{1,3} \mapsto \mathbb{C}^4
\end{align}

\section{The Form of the Dirac Equation}
The Dirac equation is given by
\begin{align}
    (i \hbar \gamma^{\mu} \partial_{\mu}  - m c) \psi(x) &= 0
\end{align}
Here we use Einstein summation over the components $\mu$ of space-time where 
\begin{align}
\gamma^{\mu} : \mathbb{C}^{4\times 4}
\end{align}
and $\psi$ is a bispinor field
\begin{align}
    \psi: \mathbb{R}^{1,3} \mapsto \mathbb{C}^4
\end{align}
%\textcolor{red}{TODO: One section on Noether's theorem is added in, add more about conserved currents here.}

%\textcolor{red}{TODO: Introduce natural units}

In natural units, this equation can be rewritten as
\begin{align}
    (i\slashed{\partial} - m)\psi(x) &= 0
\end{align}
where we define
\begin{align}
    \slashed{\partial} &= \gamma^\mu \partial_\mu
\end{align}
The Dirac equation provides us a set of 4 coupled partial differential equations. Note that in particular, if we apply $(i\slashed{partial} + m$ as an operator to both sides of the Dirac equation, we obtain the Klein-Gordon equation
\begin{align}
    (i\slashed{\partial} - m)(i\slashed{\partial} + m)\psi &= 0 \\
    (\partial_\mu \partial^\mu + m^2) \psi &= 0
\end{align}
So each component of the Dirac equation satisfies the Klein-Gordon equation. 
%\textcolor{red}{TODO: Add details on spinor-spinor products and vector-spinor products}

\section{Invariances}

The Dirac equation is invariant under actions of the Lorentz group. This follows from the use of the covariant derivative $\partial_{\mu}$. %\textcolor{red}{TODO: Is there a decent proof of this?}

\section{Solving the Dirac Equation}

When solving Dirac's equation, it is useful to start from an ansatz about the solution
\begin{align}
    \psi(x) &= u(p') e^{-p \dot x}
\end{align}
Here $p, x \in \mathbb{R}^{1,3}$ are 4-vectors where $p$ is the momentum.
\begin{align}
p &= (E_p, p')    
\end{align}
%\textcolor{red}{TODO: I don't see this following?}. 
This ansatz leads to
\begin{align}
    (\gamma^{\mu} p_{\mu} - m)u(p')
\end{align}
For a particle at rest, the momentum term disappears so we have a simplified equation
%\begin{align}
%    i\gamma^0 \frac{}
%\end{align}
\section{Quantizing Dirac's Equation}

Quantizing Dirac's equation becomes a good bit trickier than quantizing the Klein-Gordon equation.


\end{document}
